\section{Introduction}
\label{sec:introduction}

\subsection{Motivation}

The production of a well-structured academic document is a multi-day, multi-tool endeavour.
A researcher must simultaneously manage a bibliographic database, write and revise prose that
cites that database consistently, produce figures that meet venue formatting constraints,
reconcile her drafted sentences with the source material she has read, and budget the
incremental cost of every LLM query she issues.
These activities are traditionally performed by separate, loosely coupled tools:
a reference manager (Zotero, Mendeley), a writing environment (\LaTeX{}, Overleaf),
a figure-generation package (Matplotlib, TikZ), and no formal cost accounting.

The consequence of this fragmentation is well-documented.
Citation hallucination — the generation of bibliographic entries that do not correspond to
real publications — afflicts between 3\% and 47\% of LLM-generated academic texts,
depending on domain and model family~\cite{alkaissi2023hallucinated}.
Figure--text inconsistency, in which a diagram is described in the body but absent from the
compiled PDF, is a structural defect that peer review catches only after significant
reviewer effort~\cite{stelmakh2021novice}.
Token cost unpredictability makes systematic LLM use economically risky for individual
researchers without institutional API budgets~\cite{chen2023frugalgpt}.

The fragmentation problem is not merely a usability concern but a correctness concern.
When a researcher uses a separate reference manager and a separate writing tool, the
consistency between the bibliography database and the in-text citations depends on a
manual synchronisation step that is error-prone.
A reference can be cited in the text without having been added to the database
(producing a LaTeX ``undefined citation'' warning that may be overlooked), or
a reference can be added to the database without being cited (producing a spurious
entry in the bibliography).
Automated document generation systems that use separate LLM calls for reference
selection and for prose drafting face the same consistency problem at a larger scale:
a reference selected in one call may not be cited in another call because the two calls
share no common state.

\sys{} addresses all three failure modes simultaneously, within a single session, through
a formally specified system design that makes the relevant guarantees provable rather than
empirical.
The key insight is that the three failure modes correspond to three formal properties —
citation integrity, visual coherence, and cost monotonicity — that can be enforced by
three independent formal components whose designs are mutually compatible.

\subsection{Problem Statement}

Let $\mathcal{D}$ denote the space of well-formed \LaTeX{} documents — documents that
compile without error under a standard \LaTeX{} toolchain including \texttt{pdflatex}
and \texttt{biber}.
Given a research topic $q\in\mathcal{Q}$ and a session budget $\beta\in\mathbb{R}_{>0}$,
the document-generation problem is to compute a function
\[
  f:\mathcal{Q}\times\mathbb{R}_{>0}\;\longrightarrow\;\mathcal{D}
\]
such that the output $d=f(q,\beta)$ satisfies four simultaneous requirements:

\begin{enumerate}[label=\textbf{R\arabic*.},leftmargin=2.4em]
  \item \textbf{Compilability.} $d$ compiles without error under a standard \LaTeX{} toolchain.
  \item \textbf{Citation integrity.} Every bibliographic entry in $d$ is resolvable to a live
        scholarly source at the time of generation.
  \item \textbf{Visual coherence.} Every figure referenced in the body of $d$ is present in $d$
        and compiles independently.
  \item \textbf{Budget adherence.} The total token and image cost of producing $d$ does not
        exceed $\beta$.
\end{enumerate}

These requirements are not in general jointly achievable by a greedy single-pass LLM
query, because satisfying R2 requires external retrieval and verification that adds latency,
satisfying R3 requires a compilation oracle that must be invoked after generation, and
satisfying R4 requires a cost model that must be computed incrementally as generation proceeds.
A principled solution therefore requires a \emph{staged architecture} in which each stage
enforces one or more of the requirements and the inter-stage contracts are formally specified.

A further constraint, motivated by the academic publication context, is that the
generated document should be of sufficient quality that it can serve as a first draft
suitable for human revision, rather than requiring complete rewriting.
This quality constraint is not captured by R1--R4 alone; we formalise it through the
quality metric $Q(s)$ introduced in Section~\ref{sec:experimental}.

\subsection{Approach Overview}

\sys{} operationalises the function $f$ as a pipeline algebra $\Pi$ (Definition~\ref{def:pipeline})
whose stages are ordered morphisms in the category of document-state transformations.
The pipeline proceeds as follows:

\begin{description}[leftmargin=2.4em, style=nextline]
  \item[$F_1$ — Plan.]
    Given topic $q$, $F_1$ produces a structured outline $o$ that partitions the document
    into sections, assigns a target word count and figure quota to each section, and
    selects a primary language from the Language Partition $\mathcal{L}_{25}$
    (Definition~\ref{def:language}).
    The outline is stored as a typed record that subsequent stages can query.
  \item[$F_2$ — Extract and Verify.]
    $F_2$ retrieves candidate references from live scholarly endpoints (Semantic Scholar,
    CrossRef, and institutional repositories), scores each candidate against the
    \cgrag{} gate $\Gamma$ (Definition~\ref{def:cgrag}), and admits only those that
    receive the \texttt{enhanced} classification.
    Candidates assigned \texttt{degraded} are discarded before any draft text references them,
    preventing citation hallucination at the source.
  \item[$F_{2.5}$ — Reconcile.]
    The LCS Reconciliation Operator $\Phi_\theta$ (Definition~\ref{def:lcs}) aligns the
    extracted reference texts with the planned section outlines, ensuring that the
    drafting stage begins with a semantically coherent context window in which reference
    sentences and planned prose sentences share key vocabulary.
  \item[$F_3$ — Draft.]
    Given the verified reference set and the reconciled outline, $F_3$ generates section
    prose via a sequence of \texttt{chatCompletionLong} calls.
    Each call is bounded by the remaining token budget maintained by the Billing Operator
    $\mathcal{B}$ (Definition~\ref{def:billing}), ensuring that budget adherence (R4)
    is maintained throughout the drafting process.
  \item[$F_4$ — Assemble.]
    $F_4$ wraps each section in the appropriate \LaTeX{} structural commands, inserts figure
    \texttt{\textbackslash input} directives, and composes the document body.
    The order-preservation guarantee of \texttt{concurrentMap} (Theorem~\ref{thm:concurrent})
    ensures that the assembled section order matches the planned outline regardless of
    how many sections were drafted in parallel.
  \item[$F_5$ — Validate.]
    $F_5$ invokes the \LaTeX{} compiler as a formal correctness oracle.
    A non-zero exit code triggers the Typed Visual Grammar repair cycle
    (\tvg{}, Definition~\ref{def:tvg}) for any figure that failed to compile.
    The repair convergence theorem (Theorem~\ref{thm:repair}) bounds the failure
    probability, establishing compilability (R1) as a probabilistic guarantee.
\end{description}

Figure~\ref{fig:pipeline} illustrates the category-theoretic composition of these stages.

\subsection{Claims of Novelty}

The following contributions are, to our knowledge, novel:

\begin{itemize}
  \item A \textbf{formal pipeline algebra} for document generation that admits
        category-theoretic reasoning about composition and invariant preservation.
        Prior work on LLM pipelines (LangChain~\cite{chase2022langchain},
        LlamaIndex~\cite{liu2022llamaindex}) provides compositional abstractions
        without formal type systems or invariant preservation guarantees.
  \item The \textbf{\cgrag{} gate}, a multi-criteria lattice filter for reference
        verification that simultaneously gates on confidence score, citation count,
        and keyword relevance.
        Prior RAG systems gate on at most one criterion at a time.
  \item The \textbf{\tvg{} compile-repair cycle} with a closed-form convergence bound
        that depends only on per-attempt success probability and attempt count.
        Prior figure-generation systems (PlotCoder~\cite{chen2021plotcoder}) use
        compile feedback but provide no convergence guarantees.
  \item The \textbf{LCS Reconciliation Operator} as a principled mechanism for
        style-preserving source alignment, in contrast to embedding-similarity approaches
        that discard word-order information.
  \item A \textbf{unified formal model} $\mathcal{P}$ that places all components under
        a common specification, enabling joint reasoning about correctness, cost,
        and ordering — the first such unified model in the academic document
        generation literature, to our knowledge.
\end{itemize}

\subsection{Report Structure}

Section~\ref{sec:related} surveys related work across automated writing assistance,
retrieval-augmented generation, and formal pipeline specification.
Section~\ref{sec:architecture} presents the full formal specification of \sys{},
including all ten Definitions and the five system invariants.
Section~\ref{sec:contributions} derives the two main theorems and two propositions.
Section~\ref{sec:experimental} describes the evaluation methodology.
Section~\ref{sec:results} presents quantitative results across four evaluation objectives.
Section~\ref{sec:discussion} discusses the interpretation of results, the scope
and limitations of the formal model, and four future research directions.
Section~\ref{sec:conclusion} concludes.
